\section{A tame deformation and its Selmer length}\label{sec:tame}

We retain the fields and the contradiction assumption of
\cref{sec:reduction}.  In particular, the conductor-one Heegner trace
$y\in A_0(K)$ is torsion.  Put $T=T_pA_0$ and
$V=T\otimes_{\Z_p}\Q_p$.  Initially all cohomology in this section has
$\Q_p$ coefficients; extension to the fixed finite coefficient field
$k$ is made by tensor product.  Enlarge the finite, conjugation-stable
set $S$ to contain the places over $p$, the bad places of $A_0$, and all
the fixed conductors used below.  It may be enlarged by a further
\emph{fixed} finite set at any stage.  In particular, the auxiliary
finite characters and tests of Section~\ref{sec:content} are chosen
independently of the moving primes, and their conductors are included
before the final prime sequence is realized.  All the constructions
below allow this fixed enlargement.  Write $G_{K,S}$ for the Galois
group unramified outside $S$.

We need two conclusions for the same tame deformation.  The
one-sided local condition, full at one $p$-adic place and zero at
its conjugate, will give a cohomology module of length zero or at
most two.  Under the torsion assumption on $y$, the first one or
three coefficients of the corresponding weighted Heegner logarithms
will tend to zero, respectively.  These are
\cref{prop:green-length,prop:heegner-log}.  The logarithm estimate
supplies the constant-term estimate of
\cref{prop:constant-term-small}; the length bound will contradict
the extensions extracted in Section~\ref{sec:extensions}.

\subsection{The undeformed Selmer spaces at the two places above \texorpdfstring{$p$}{p}}

The two places $v,v^c$ of $K$ above $p$ have completion $\Q_p$.
The alternating Weil pairing
\[
 e:V\times V\longrightarrow\Q_p(1)
\]
and the local invariant give a symmetric perfect pairing on
$H^1(K_w,V)$: the sign from interchanging two degree-one cochains
cancels the sign of $e$.  The local Kummer line
$H^1_f(K_w,V)$ is its own orthogonal complement.  Thus, for $w\mid p$,
this is a hyperbolic plane.  These statements include arbitrary
reduction at $p$: local torsion is finite, the rational Kummer space
has dimension one, and the local Euler characteristic gives
$\dim H^1(K_w,V)=2$.  At a finite place $w\nmid p$ we instead have
\begin{equation}\label{eq:away-p-vanishing}
 H^0(K_w,V)=H^2(K_w,V)=H^1(K_w,V)=0.
\end{equation}
Indeed $H^0=0$ by finiteness of local torsion, $H^2=0$ by local
duality, and the Euler characteristic away from $p$ is zero.
We use the local and global duality sequences in
\cite[Chapter~I, Sections~2--4]{Milne1986} and
\cite[Section~8.6]{NSW2008}.

For the moment let $W$ denote either $V$ or its twist $V^{DD'}$ over
$K$.  Its compact Selmer space has dimension one by
\cref{prop:arithmetic-reduction}.  Let $H^1_{\mathrm{rel}}(K,W)$ have full
conditions at $v,v^c$, and the usual unramified conditions outside
$S$; conditions at the other places of $S$ make no difference by
\cref{eq:away-p-vanishing}.  Let $H^1_{\Gr}(K,W)$ have full condition
at $v$ and zero condition at $v^c$, and let
$H^1_{\overline{\Gr}}(K,W)$ have the opposite conditions.
These are the two one-sided Greenberg conditions.

\begin{lemma}\label{lem:base-green}
If the compact Selmer line of $W/K$ has nonzero localization, then
$H^1_{\Gr}(K,W)=H^1_{\overline{\Gr}}(K,W)=0$.
The only remaining possibility is the following \emph{strict case}
in the $V$ summand: $r=1$, and the compact Selmer space is generated
by a class $a$ such that $a^c=a$ and every localization of $a$ is zero.
In that case
\begin{equation}\label{eq:base-green-bases}
 H^1_{\Gr}(K,V)=\langle a,h\rangle,
 \qquad
 H^1_{\overline{\Gr}}(K,V)=\langle a,h^c\rangle,
\end{equation}
where $h$ has nonzero localization at $v$ and zero localization at
$v^c$.  The $V^{DD'}$ summand always has zero Greenberg space.
\end{lemma}

\begin{proof}
Let $F_w=H^1_f(K_w,W)$ and let $Q_w=H^1(K_w,W)/F_w$.
Poitou--Tate duality identifies the image of
\[
 H^1_{\mathrm{rel}}(K,W)\longrightarrow Q_v\oplus Q_{v^c}
\]
with the annihilator of the localization of the compact Selmer line
in $F_v\oplus F_{v^c}$.  If that localization is nonzero, both of its
coordinates are nonzero: the one-dimensional global space is stable
under $c$, which interchanges the two places.  Its annihilator is
therefore a line whose two coordinates are both nonzero for every
nonzero vector.  A class satisfying a zero condition at one place
has zero singular image, hence belongs to the compact Selmer line;
the zero condition then kills that line as well.

Suppose the compact Selmer line is everywhere locally zero.  The
same duality sequence makes the singular image all of
$Q_v\oplus Q_{v^c}$.  The kernel of localization is precisely the
compact Selmer line.  Consequently the relaxed local image has
dimension two and maps isomorphically to the singular plane.
In each local hyperbolic plane choose the isotropic complement $E_w$
of $F_w$, choosing the two complements compatibly with $c$.  The
relaxed local image is the graph of a map
\[
 E_v\oplus E_{v^c}\longrightarrow F_v\oplus F_{v^c}.
\]
Global reciprocity makes the graph isotropic.  In mutually dual
bases its matrix is therefore skew-symmetric, hence of the form
$\left(\begin{smallmatrix}0&u\\-u&0\end{smallmatrix}\right)$.
Conjugation interchanges both pairs of basis vectors and preserves
the local pairings.  Invariance of the graph therefore changes $u$
to $-u$, so $u=0$.  Lifting the first coordinate line gives $h$;
conjugating gives $h^c$ and proves \cref{eq:base-green-bases}.

The auxiliary summand has a nontorsion global point and nonzero
localization, by \cref{prop:arithmetic-reduction}.  If $r=0$, the unique
Selmer line of $V/K$ likewise comes from a nontorsion point on the
quadratic twist $A_0^D/\Q$, so it cannot be strict.  When $r=1$, the
twist by $D$ has zero Selmer space and the line comes from
$H^1_f(\Q,V)$, giving $a^c=a$.
\end{proof}

\subsection{Primes with a controlled tame logarithm}

In the prime construction and cohomological calculations of this section,
Galois Frobenius elements are arithmetic Frobenius elements.
The following construction gives a tame direction that is trivial at
the fixed places and has uniform local bounds.  In the strict case
we will refine its limiting Frobenius element after determining the
obstruction to first-order lifting.

\begin{proposition}\label{prop:tame-primes}
There are rational primes $q_i\notin S$, split completely in $L$,
distinguished places $\mathfrak q_i$ of $K$ above them, cyclic
$p$-quotients $\Delta_i$ of the ring class group of conductor $q_i$,
and surjections
\[
 \ell_i:G_K\longrightarrow\Delta_i\simeq\Z/p^{n_i}\Z,
 \qquad n_i\longrightarrow\infty,
\]
with the following properties.
\begin{enumerate}[label=\textup{(\roman*)}]
\item $\ell_i$ is unramified outside the pair above $q_i$,
its restriction to inertia at $\mathfrak q_i$ is surjective,
$\ell_i^c=-\ell_i$, and $\ell_i|_{G_{K_s}}=0$ for every $s\in S$.
\item $v_p(q_i-1)-n_i$ is bounded.  At either place above $p$, the
full ring class field of conductor $q_i$ is unramified and its local
degree has bounded $p$-part.
\item Frobenius at $\mathfrak q_i$, on all fixed data unramified
outside $S$, tends to an element $\gamma\in G_L$ such that
\[
 \eps(\gamma)=1,\qquad
 \det(V(\gamma)-1)\ne0,\qquad
 \cC(\gamma)\text{ is not a root of unity}.
\]
Moreover $\gamma$ fixes all $p$-power radicals of a fixed,
conjugation-stable finite list in $K^\times$ spanning
$\cO_{K,S}^{\times}\otimes\Q_p$ under the Kummer map.
\end{enumerate}
\end{proposition}

\begin{proof}
Choose representatives $\mathfrak a_j$ for generators of the ideal
class group, prime to $S$.  For their orders, their multiplication
relations, and the expressions of the primes in $S$ in these
generators, choose principal generators.  Adjoin generators of the
finite-rank $S$-unit group and their conjugates to this finite list,
denoted $\alpha_1,\ldots,\alpha_d$.  Fixed class-group orders and
roots of unity will cost only a fixed index.

For a split prime $q$, the kernel of the ring class group map to the
ordinary ideal class group is the quotient of
$(\cO_K/q\cO_K)^\times$ by diagonal residue units and global units.
Its $p$-part is cyclic, of order $p^{v_p(q-1)}$ up to a fixed factor.
Require the $\alpha_j$ to have zero residue logarithm modulo
$p^{n_i}$ at the pair over $q_i$.  This is ensured by splitting in
their $p^{n_i}$-radical extensions, together with the corresponding
roots of unity.  The ring class presentation now allows us to set
the logarithms of all $\mathfrak a_j$ equal to zero and to take the
surjective logarithm on the residue quotient.  Each relation is
respected, and the expressions chosen for the primes of $S$ give
$\ell_i|_{G_{K_s}}=0$.  Conjugation inverts the residue quotient and
the ring class action, giving $\ell_i^c=-\ell_i$.

We verify that imposing the full radical condition leaves room for
$\gamma$.  Over the cyclotomic tower the Kummer Galois group is a
closed subgroup of $\Z_p(1)^d$.  The joint image of $V$ and $\cC$
is a $p$-adic Lie group.  A common Lie quotient with this translation
group is abelian, so it kills the derived Lie algebra of the image
of $V$.  On the remaining torus Lie algebra, conjugation by an open
cyclotomic subgroup is trivial; on the Kummer translations it is
multiplication by the cyclotomic character.  Any equivariant common
quotient is consequently zero.  This argument can be made after a
fixed finite extension, which does not change Lie dimensions.
Thus killing the radicals removes no Lie dimension from the joint
image over the cyclotomic tower.

For a non-CM elliptic curve the determinant-one image has Lie
algebra $\mathfrak{sl}_2$, by the open image theorem
\cite{Serre1972}; for a CM curve its identity component contains
the norm-one torus, by the CM description of the Tate-module
character \cite{Serre1968}.  In either
case the locus $\det(V(g)-1)\ne0$ is nonempty and open in every
sufficiently small identity neighborhood.  The infinity type of
$\cC$ makes its restriction over the cyclotomic tower nonconstant;
its torsion locus is a proper closed subset in a sufficiently small
identity neighborhood.  The two complements intersect.  We may
also impose $g\in G_L$ and any fixed open-kernel conditions.  This
gives $\gamma$ fixing the full radicals and the cyclotomic tower
with the required properties.

It remains to control $v_p(q_i-1)$, rather than merely make it
large.  The cyclotomic--radical image is a closed subgroup of
\[
 \Z_p(1)^d\rtimes\Z_p^\times
\]
with open cyclotomic projection.  Choose an element in a pro-$p$
subgroup whose cyclotomic projection is nontorsion.  Its $p^m$-th
powers tend to the identity; their translation coordinates are
divisible by $p^{m-C_0}$, and the valuation of their cyclotomic
coordinate minus one is $m+C_1$, for fixed $C_0,C_1$.  Multiplying
$\gamma$ by such a power with $m=n_i+C_2$ therefore leaves all
radicals trivial modulo the requested precision and makes the
cyclotomic valuation exactly $n_i+C$, for a fixed $C$.  Chebotarev
applied to a growing sequence of finite quotients realizes these
conditions and increasingly accurate approximations on all fixed
data.  Choosing a place above each prime selects the indicated
Frobenius representative.  The primes can be taken outside any
fixed finite set.

Finally the $p$-part of the order of the full ring class group is
$p^{v_p(q_i-1)+O(1)}$.  The kernel of its map to $\Delta_i$ therefore
has bounded $p$-part.  The Frobenius of every $s\in S$ belongs to this
kernel.  At $p$ the extension is unramified, so the $p$-part of the
local degree, which is the $p$-part of the Frobenius order, is
bounded.  This proves the proposition.
\end{proof}

\subsection{Admissible cochains and division by the parameter}

Fix a nonprincipal ultrafilter $\cU$ on the prime indices.  A
statement ``eventually'' in the construction below means on a set
in $\cU$.  For every fixed $b\ge1$, put
\[
 R_b=k[t]/(t^b).
\]
At finite stages the substitution of $(1+t)^{\ell_i}$ is defined
only to finite $p$-adic precision.  There is a constant $c_b$ such that
\begin{equation}\label{eq:binomial-loss}
 u\equiv u'\pmod {p^n}\quad\Longrightarrow\quad
 \binom{u}{j}\equiv\binom{u'}{j}\pmod {p^{n-c_b}}
 \quad(0\le j<b).
\end{equation}
This follows by writing each binomial coefficient as a polynomial
with denominator $j!$.  Thus truncated substitution of
$(1+t)^{\ell_i}$ defines an action modulo $p^{n_i-c_b}$ and sends
integral group-ring elements of $\Delta_i$ to coefficients with a
fixed denominator.  No division by $|\Delta_i|$ is involved.

For a sequence $(g_i)$, let $\ell((g_i))$ be the $p$-adic ultralimit
of $\ell_i(g_i)\in\Z/p^{n_i}\Z$: each fixed residue modulo $p^m$ is
well-defined eventually and has a unique ultralimit.  This is an
additive character on the group of sequences of elements of
$G_{K,Sq_i}$, modulo equality on a set in $\cU$.  We can therefore
define
\[
 \psi=(1+t)^\ell,\qquad
 M_b^{\eta}=V\otimes_{\Q_p}R_b(\psi^{\eta}),
 \qquad \eta\in\{1,-1\}.
\]
Constant representations have their usual compact $p$-adic ultralimit.

\begin{definition}\label{def:admissible}
An \emph{admissible cochain} with coefficients in $M_b^{\eta}$ is a
function on sequences of Galois elements which is the coefficientwise
ultralimit of continuous finite-stage cochains, in the fixed basis
of $V\otimes R_b$, with one denominator $p^C$ independent of $i$.
Finite-stage cochain identities are required modulo precisions
$A_i\longrightarrow\infty$ along $\cU$, with
$A_i\le n_i-c_b$.  Equivalently, the limiting identity holds on
\emph{every} sequence of arguments.  Cohomology is the cohomology
of this cochain complex, so coboundaries are also required to be
admissible.  Local complexes are defined by the chosen embeddings
of decomposition groups.  The same definition applies over $L$ by
restriction and inverse image of the finite-stage Galois groups.
\end{definition}

The equivalence in the definition uses boundedness: if an identity
failed uniformly modulo $p^m$ on a set in $\cU$, choosing one
violating tuple at every such index would violate the identity on
that sequence.  A diagonal choice of precisions then gives
$A_i\to\infty$.  This observation also applies when replacing two
cochain systems with the same limit.  In particular, pointwise
agreement only on fixed Galois elements would not suffice.

\begin{lemma}\label{lem:admissible-specialization}
For constant coefficients $V$, admissible $H^1$ is precisely
$H^1(G_{K,S},V)$, and likewise over $L$.  For $\Q_p(1)$ the same
assertion holds when unramifiedness at the moving primes is imposed.
The fixed local conditions specialize to the usual rational local
conditions.  Coefficient reduction and multiplication by $t$ give
short exact sequences of admissible cochain complexes
\begin{equation}\label{eq:coefficient-sequence}
 0\longrightarrow C^\bullet(M_{b-1}^{\eta})
 \xrightarrow{\ t\ } C^\bullet(M_b^{\eta})
 \longrightarrow C^\bullet(V\otimes k)\longrightarrow0.
\end{equation}
Zero local conditions at $p$ and local Kummer conditions are
preserved by the resulting division of a class with zero
specialization by $t$.
\end{lemma}

\begin{proof}
Choose a stable lattice in $V$.  A bounded finite-stage cocycle
modulo $p^{A_i}$ has trivial wild inertia at $q_i$.  On tame
$p$-inertia the Frobenius relation implies
\[
 (V(\Frob_i)-q_i)z_i(\tau)=0\pmod {p^{A_i-C}}.
\]
The inverses of $V(\Frob_i)-q_i$ have a fixed denominator: these
matrices converge to the invertible matrix $V(\gamma)-1$.
The same is true at the conjugate prime and over $L$.  After one
fixed loss of precision, the cocycle kills the normal subgroup
generated by moving inertia, hence factors through $G_{K,S}$.
For each fixed finite coefficient module, there are only finitely
many continuous cocycles on $G_{K,S}$: cohomological finiteness
follows from the finiteness of extensions with fixed ramification
and bounded degree, and the group of coboundaries is finite.
Taking the compatible ultralimits of these finite sets yields an
actual continuous integral cocycle; then undo the fixed scaling.
For cyclotomic coefficients the proof is the same once moving
inertia has been required to vanish.  Equality of classes is
compatible with this procedure because bounded zero-cochains
have compact ultralimits.  It also shows that a representative of
a given ordinary class can be replaced by a chosen continuous
representative, up to a bounded coboundary and uniformly vanishing
errors.

Exactness of \cref{eq:coefficient-sequence} is coefficientwise:
lift the finitely many coordinates continuously, and divide a
zero constant coefficient by $t$ by shifting coordinates.  These
operations do not enlarge a uniform denominator except for the
fixed truncation loss in \cref{eq:binomial-loss}.  The character is
exactly trivial at the fixed places in $S$, so there
\[
 H^1(K_s,M_b^{\eta})=H^1(K_s,V\otimes k)\otimes_k R_b,
 \qquad
 H^1_f(K_s,M_b^{\eta})=H^1_f(K_s,V\otimes k)\otimes_k R_b.
\]
These are fixed linear subspaces.  A strictly zero local class
can be made a zero local cocycle with bounded corrections: since
$H^0(K_s,V)=0$, finitely many elements $g_1,\ldots,g_u$ make
\[
 V\longrightarrow V^u,\qquad
 A\longmapsto((V(g_j)-1)A)_j
\]
injective.  A fixed left inverse bounds the denominator of every
such correction.  Therefore division preserves the strict and
Kummer conditions, both in rational cohomology and in their
finite-precision representatives.
\end{proof}

Over $L$ we impose full local conditions at the places of
$\Sigma_p$, zero conditions at their conjugates, and relaxed
conditions at the other places of $S$.  This is the notation
$\GrH(L,M_b^{\eta})$ used below, always with admissible cohomology.
Restriction, the two idempotents for $\Gal(L/K)$, and corestriction
decompose it into the corresponding spaces for $V$ and $V^{DD'}$
over $K$.  Their denominators are fixed (in fact $2$ is a unit).
Thus all length computations may be made over $K$.

\subsection{The first-order obstruction}

We now work in the strict case of \cref{lem:base-green}.  Choose
fixed continuous cocycles for $a,h,h^c$, taking the representatives
equivariantly under $c$ by averaging.  Cup products below contract
with $e$, and $\delta$ denotes the cochain differential.

Let $x$ be a fixed cocycle representing a class in
$H^1_{\Gr}(K,V)$.  A first-order lift has the form $x+t x_i'$.
Comparing the coefficients of $t$ in its cocycle
equation gives
\begin{equation}\label{eq:linear-lift}
 \delta x_i'=-\eta\ell_i\cup x
\end{equation}
modulo precision tending to infinity.  To obstruct a lift with the
Greenberg local conditions, we pair this equation with a class
having the opposite local conditions.  The proof below justifies
the bounded corrections needed to use the fixed representative $x$.

For $d\in H^1_{\overline{\Gr}}(K,V)$, the
product $x\cup d$ vanishes at every place: at $p$ one of the two
classes is zero, and away from $p$ use
\cref{eq:away-p-vanishing}.  The injection of global degree-two
cyclotomic cohomology into its localizations (Brauer reciprocity)
gives a continuous cochain $u$ on $G_{K,S}$ satisfying
\begin{equation}\label{eq:cup-primitive}
 \delta u=x\cup d.
\end{equation}
Choose primitives for the finitely many pairs in the two fixed
bases, and thereafter use bilinear extension.  Each has compact
image, hence all these choices have one fixed denominator.

For an element $\gamma$ as in \cref{prop:tame-primes}, define
\begin{align}
 z_\gamma(x,d;u)
 &=u(\gamma)-e\bigl((V(\gamma)-1)^{-1}x(\gamma),d(\gamma)\bigr),
 \label{eq:corrected-primitive}\\
 H_\gamma(x,d)
 &=z_\gamma(x,d;u)-z_\gamma(x^c,d^c;u^c).
 \label{eq:obstruction-pairing}
\end{align}
The chosen trivialization of $\Q_p(1)$ makes these scalar
expressions; changing it multiplies them by one nonzero scalar.

\begin{lemma}\label{lem:obstruction-necessary}
The expression $H_\gamma$ is independent of the primitive and
bilinear.  If a nonzero class $x\in H^1_{\Gr}(K,V\otimes k)$ is the
specialization of an admissible class in
$\GrH(K,M_2^{\eta})$, then
\begin{equation}\label{eq:obstruction-vanishing}
 H_\gamma(x,d)=0
 \quad\text{for every }d\in H^1_{\overline{\Gr}}(K,V\otimes k).
\end{equation}
The same necessary condition holds for either $\eta$.
\end{lemma}

\begin{proof}
We give the finite-precision argument, so that no continuity
assertion about the varying Galois groups is implicit.  For
$\eta=1$, write a first-order lift as $x+t x_i'$ satisfying
\cref{eq:linear-lift}.  Indeed,
\cref{lem:admissible-specialization} permits subtraction of one
bounded global coboundary and replacement of its constant
coefficient by the chosen fixed cocycle $x$.  Equality on all
sequences makes the error uniformly small.  Its linear coefficient
remains bounded.  The cochain
\begin{equation}\label{eq:global-two-cocycle}
 x_i'\cup d-\ell_i\cup u
\end{equation}
is consequently a two-cocycle modulo the reduced precision.
Its local invariants at $S$ vanish to increasing precision.  At
$p$ this follows from opposite strict conditions and $\ell_i=0$;
away from $p$ it follows from the vanishing of rational $H^1$ and
the bounded exponents of the fixed integral local cohomology
groups.  At places outside $S q_i$ it is unramified and has zero
invariant.

At $\mathfrak q_i$, the fixed cocycles $x,d$ are unramified and
are coboundaries, say $x=\delta A_i$, $d=\delta B_i$, where
\[
 A_i=(V(\Frob_i)-1)^{-1}x(\Frob_i),\qquad
 B_i=(V(\Frob_i)-1)^{-1}d(\Frob_i).
\]
These vectors have bounded denominators, because the inverse
matrices converge to $(V(\gamma)-1)^{-1}$.  By
\cref{eq:linear-lift}, $x_i'-\ell_i\cup A_i$ is a cocycle.
Its cup product with $d$ is a coboundary.  The invariant of
\cref{eq:global-two-cocycle} therefore equals minus the invariant
of
\begin{equation}\label{eq:tame-corrected-cup}
 \ell_i\cup(u-A_i\cup d).
\end{equation}
The second factor is an unramified cyclotomic one-cocycle locally,
since its differential is $x\cup d-(\delta A_i)\cup d=0$.
The tame local pairing evaluates this unramified factor at
Frobenius and pairs it with the inertia logarithm.  With compatible
local invariant conventions, the multiplier at stage $i$ is a unit
$\nu_i$ modulo the working precision, because the logarithm on tame
inertia is surjective.  At the conjugate prime the same calculation
uses $x^c,d^c,u^c$ and $\ell_i^c=-\ell_i$; functoriality of the local
invariant gives the same multiplier $\nu_i$.  Choose unit lifts to
$\Z_p$ and let $\nu\in\Z_p^\times$ be their compact ultralimit.
The global sum of invariants then gives $\nu H_\gamma(x,d)=0$,
hence \cref{eq:obstruction-vanishing}.  This calculation never
divides by $q_i-1$.  We use arithmetic Frobenius
in this calculation; replacing it throughout by geometric
Frobenius changes the conventions compatibly and not the zero
condition.  For $\eta=-1$, replace $\ell_i$ by $-\ell_i$.

Here is the precise bound on all reductions of precision.  Let
$c_{\mathrm{fix}}$ bound the fixed cocycles and their finitely many
primitives, and let $c_{\mathrm{lift}}$ bound the particular
admissible lift.  Let $c_\gamma$ clear both inverses
$V(\Frob_i)-1$ and $V(\Frob_i)-q_i$ eventually.  Let $c_S$ bound
the fixed left inverses detecting local $H^0=0$ and the exponents
of integral $H^1,H^2$ at the fixed places away from $p$.  Finally
use the binomial constant $c_b$ of \cref{eq:binomial-loss}.
All displayed corrections use finitely many sums and products of
these objects.  A constant
$C=C(c_{\mathrm{fix}},c_{\mathrm{lift}},c_\gamma,c_S,c_b)$ therefore
bounds every loss.  If the original cochain identities hold to
precision $A_i$ and their replacement accuracies are $B_i$, make
the calculation to any precision
\begin{equation}\label{eq:precision-choice}
 N_i\le\min\{n_i,A_i,B_i\}-C,\qquad N_i\longrightarrow\infty.
\end{equation}
The tame pairing above is defined at this precision since
$N_i\le n_i$.  The constants may depend on a fixed lift and on
$b$, but never on $i,q_i,n_i$.  Only finitely many fixed primitives
were chosen; no bounded contracting homotopy on all cochains is
being asserted.

For independence of $u$, the difference of two primitives is in
\[
 H^1(G_{K,S},\Q_p(1))
   =\cO_{K,S}^{\times}\otimes_{\Z}\Q_p.
\]
The equality is the Kummer sequence and finiteness of the
$S$-class group.  The radical list in \cref{prop:tame-primes}
spans this space, including its conjugate, up to a fixed index.
Every such Kummer class evaluates to zero at $\gamma$; a
coboundary also evaluates to zero because $\eps(\gamma)=1$.
Changing $u$ therefore changes neither term of
\cref{eq:obstruction-pairing}.  Bilinearity follows from the fixed
bilinear choices of primitives.
\end{proof}

\begin{lemma}\label{lem:nonzero-obstruction}
The element $\gamma$ in \cref{prop:tame-primes} can be chosen so
that, in the bases of \cref{eq:base-green-bases}, the obstruction
matrix is
\begin{equation}\label{eq:obstruction-matrix}
 \begin{pmatrix}0&\kappa\\-\kappa&0\end{pmatrix},
 \qquad \kappa\in\Q_p^\times.
\end{equation}
Thus no nonzero base Greenberg class lifts to order two.
\end{lemma}

\begin{proof}
First $H_\gamma(x,d)=H_\gamma(d,x)$ whenever the pairs under
consideration have locally zero cup product.  To check the sign
explicitly, put $q(g)=e(x(g),d(g))$.  The cocycle identities give
\[
 \delta q=d\cup x-x\cup d,
\]
so $u+q$ is a primitive for the switched cup product.  At a place
where $x=\delta A$ and $d=\delta B$, the difference between the two
corrected evaluations is
$(\eps(\gamma)-1)e(A,B)=0$.  This proves the symmetry for the
expression in \cref{eq:corrected-primitive}, and hence for $H$.
Its definition also gives
$H_\gamma(x^c,d^c)=-H_\gamma(x,d)$.  Since $a^c=a$ and $p>2$,
these identities give both diagonal zeros in
\cref{eq:obstruction-matrix} and the opposite signs of its
off-diagonal entries.

It remains to make $\kappa=H_\gamma(a,h^c)$ nonzero.  Set
$d^-=h^c-h$.  It is anti-invariant, and its cup product with $a$
is everywhere locally zero.  Average a primitive to be
anti-invariant as well.  Then
\[
 H_\gamma(a,d^-)=2z_\gamma(a,d^-;u),\qquad
 H_\gamma(a,d^-)=2H_\gamma(a,h^c).
\]
It is enough to make $z_\gamma(a,d^-;u)$ nonzero.

Let $G=G_{K,S}$ and $N_V=\ker(G\to\GL(V))$.  The classes $a,d^-$
are linearly independent: the former is nonzero and invariant
under $c$, whereas the latter is anti-invariant with nonzero local
image.  The image of $G$ on $V$ contains a central scalar different
from one.  The usual central-scalar argument gives
$H^1(\im(G\to\GL(V)),V)=0$: for a cocycle $f$ and a central
scalar $z\ne1$, the cocycle identity is
$(z-1)f(g)=(g-1)f(z)$, so $f$ is a coboundary.
Inflation--restriction thus makes restriction to
$N_V$ injective on $H^1(G,V)$.  On $N_V$, both cocycles are additive
homomorphisms.  Their joint image spans $V\oplus V$: its span is a
$G$-submodule, and absolute irreducibility with
$\End_G(V)=\Q_p$ says that a proper submodule projecting onto each
factor would impose a scalar linear relation between the two
restrictions.  That contradicts the preceding injectivity.
The absolute irreducibility here follows from open image in the
non-CM case; in the CM case it follows because neither $K$ nor $L$
contains the CM field.  Passing to the finite-index subgroup over
$L$ leaves the span unchanged.

The Weil pairing implies $\eps=1$ on $N_V$.  On this group,
\cref{eq:cup-primitive} reads
\[
 u(gg')=u(g)+u(g')-e(a(g),d^-(g')).
\]
Consequently
\begin{equation}\label{eq:heisenberg-commutator}
 u([g,g'])=-e(a(g),d^-(g'))+e(a(g'),d^-(g)).
\end{equation}
The alternating form on $V\oplus V$ on the right is nonzero and
nondegenerate.  Since the joint translations span $V\oplus V$,
some $g,g'\in N_V\cap G_L$ have a nonzero right-hand side.
Write $\sigma=[g,g']$.  Both translations vanish on $\sigma$, but
$u(\sigma)\ne0$.

Every Kummer character in our radical list is an additive
homomorphism on $N_V$, because $\eps=1$ there.  It therefore
vanishes on $\sigma$, so $\sigma$ fixes all the radicals.  A
one-dimensional character such as $\cC$ also kills $\sigma$.
Replacing $\gamma$ by $\gamma\sigma^m$, for an integer $m$, thus
preserves its $V$, $\cC$, cyclotomic, and radical values.  The two
translations entering the correction in
\cref{eq:corrected-primitive} do not change, whereas its primitive
value changes by $m u(\sigma)$.  Some $m$ makes the result nonzero.
If further fixed open conditions are required, first replace
$\sigma$ by a suitable nonzero power; its nonzero primitive value
survives.  This constructs the required $\gamma$.  In particular
the inverse matrix $V(\gamma)-1$, and hence its denominator bound,
is unchanged by the modification.  Applying the prime construction
to this element proves \cref{eq:obstruction-matrix} for both signs
of the tame twist.
\end{proof}

In the strict case we henceforth use this refined $\gamma$ in
\cref{prop:tame-primes}, and $q_i,\ell_i$ denote the sequence
realized from it.  With the same ultrafilter $\cU$, we henceforth
form $\ell$, $\psi$, $M_b^\eta$, and the admissible global and local
cochain complexes from this reselected sequence.  The preserved
properties of the prime construction ensure that the specialization,
division, and obstruction arguments above still apply.
As stipulated at the start of the section, the
final realization is made after all fixed conductors and tests have
been chosen.  Any additional fixed open-kernel conditions can be imposed
in the construction and preserved by the commutator modification.

\subsection{The exact Artin length bound}

\begin{proposition}\label{prop:green-length}
For the primes of \cref{prop:tame-primes}, chosen in the strict case
as in \cref{lem:nonzero-obstruction}, every fixed $b\ge1$, and
either $\eta=1$ or $\eta=-1$, one has
\begin{equation}\label{eq:green-length}
 \length_{R_b}\GrH(L,V\otimes\psi^{\eta})
 \begin{cases}
 =0,&\text{outside the strict case},\\
 \le2,&\text{in the strict case}.
 \end{cases}
\end{equation}
In the strict case this module is killed by $t$ and is the image
of its two-dimensional base space under multiplication by
$t^{b-1}$.  Put
\begin{equation}\label{eq:b-star}
 b_*=1\quad\text{outside the strict case},\qquad
 b_*=3\quad\text{in the strict case}.
\end{equation}
Thus in both cases the length in \cref{eq:green-length} is
strictly smaller than $b_*$.
\end{proposition}

\begin{proof}
It suffices to work in each of the two summands over $K$.
Let $\mathcal S_b^\bullet$ be the cone, shifted by $-1$, of the
map from the admissible global cochain complex to the local
complex at the strict $p$-place.  Its degree-one cohomology is
exactly the indicated Greenberg group.  Indeed global and local
$H^0$ vanish for $V$ and, by the $t$-filtration, for every $M_b$.
The coefficient sequence in \cref{eq:coefficient-sequence} gives a
short exact sequence of these cone complexes.  The resulting
cohomology sequence therefore identifies
\begin{equation}\label{eq:selmer-kernel}
 \ker\bigl(H^1(\mathcal S_b^\bullet)
       \longrightarrow H^1(\mathcal S_1^\bullet)\bigr)
 =\im\bigl(t:H^1(\mathcal S_{b-1}^\bullet)
       \longrightarrow H^1(\mathcal S_b^\bullet)\bigr).
\end{equation}
There is no flatness assumption on these cohomology groups.

If the base group is zero, induction in \cref{eq:selmer-kernel}
gives zero at every order.  This deals with the auxiliary summand
and with both summands outside the strict case.  In the strict
summand, a class at any order $b\ge2$ reduces to an order-two
class, whose base specialization is zero by
\cref{lem:obstruction-necessary,lem:nonzero-obstruction}.
Equation~\eqref{eq:selmer-kernel}, iterated, gives
\[
 H^1(\mathcal S_b^\bullet)
   =\im\bigl(t^{b-1}:H^1(\mathcal S_1^\bullet)
                    \longrightarrow H^1(\mathcal S_b^\bullet)\bigr).
\]
This image is killed by $t$ and has $k$-dimension at most two.
Its $R_b$-length is therefore at most two, as asserted.
\end{proof}

\subsection{A uniform local logarithm estimate}

The cohomological length bound is now proved.  Under the torsion
assumption on $y$, we next establish the weighted Heegner-logarithm
vanishing needed for the theta congruence.  The argument first gives
vanishing of local Kummer classes.  The following estimate converts
this into a valuation statement, uniformly in the unramified fields
of definition that occur at $p$.

\begin{lemma}\label{lem:uniform-local-log}
Fix an integer $d_0\ge0$.  There is a constant $c$ such that, for
every finite unramified extension $E/\Q_p$ with
$v_p([E:\Q_p])\le d_0$,
\begin{equation}\label{eq:uniform-local-log}
 \log_{A_0,p}(A_0(E))\subseteq p^{-c}\cO_E.
\end{equation}
If $P\in A_0(E)$ and the Kummer class of $p^C P$ vanishes in
$H^1(E,A_0[p^n])$, where $C,n\ge0$ are integers, then
\[
 v_p(\log_{A_0,p}P)\ge n-C-c.
\]
\end{lemma}

\begin{proof}
On a fixed sufficiently small formal subgroup, the formal logarithm
and its inverse are integral, uniformly under unramified base
extension.  The intervening formal-group quotient has bounded
$p$-exponent.  In the N\'eron special fiber the geometric component
group has fixed order, the additive part has bounded $p$-exponent,
and a torus in characteristic $p$ has no $p$-torsion.  For the
\'etale $p$-power torsion in the abelian part, let $\alpha$ be a
unit Frobenius eigenvalue.  It is not a root of unity, by its
complex absolute value.  If $[E:\Q_p]=d$, the exponent is
controlled by $v_p(\alpha^d-1)$.  Writing $\alpha$ as its finite
torsion part times a principal unit shows
\[
 v_p(\alpha^d-1)\le C_\alpha+v_p(d)
\]
whenever the left side is positive.  It is therefore bounded for
these extensions.  A fixed $p$-power, followed if necessary by an
integer prime to $p$, moves any point into the formal subgroup.
Dividing its logarithm by that multiplier proves
\cref{eq:uniform-local-log}, including at bad reduction.

For the precision assertion, the Kummer injection
\[
 A_0(E)/p^n A_0(E)\hookrightarrow H^1(E,A_0[p^n])
\]
gives $p^C P=p^n Q$ for some $Q\in A_0(E)$.  Consequently
\[
 v_p(\log_{A_0,p}P)\ge n-C-c
\]
by \cref{eq:uniform-local-log}.  Thus the comparison loses only
a fixed precision, without division by the local extension degree.
\end{proof}

\subsection{Weighted Heegner traces and their local logarithms}

Let $H_i/K$ be the full ring class field of conductor $q_i$, and
let $y_i\in A_0(H_i)$ be the corresponding Heegner point.  At the
split prime $q_i$ it is the conductor-$q_i$ translate obtained by
$n(1/q_i)$ in torus-adapted coordinates, with the opposite sign
if the embedding convention requires it.  Write
$\mathcal G_i=\Gal(H_i/K)$, and view $\ell_i$ also as its quotient
character.  The sums in the following statement are raw sums
over $\mathcal G_i$.

\begin{proposition}\label{prop:heegner-log}
At the distinguished place over $p$, and for either sign,
\begin{equation}\label{eq:heegner-log}
 \sum_{g\in\mathcal G_i}(1+t)^{\pm\ell_i(g)}
       \log_{A_0,p}(y_i^g)
       \longrightarrow0\pmod {t^{b_*}}.
\end{equation}
More precisely, every coefficient of degree less than $b_*$ has
valuation tending to $+\infty$ along $\cU$.  The coefficients
through any other fixed degree have a common lower valuation
bound.  These assertions hold in a common $p$-adic algebraic
closure, with its fixed extension of the valuation.
\end{proposition}

\begin{proof}
Choose integral Kummer cocycles for $y_i$ and corestrict with the
tautological group-ring coefficients.  Corestriction is a sum:
its terms have integral coefficients regardless
of the size of $\mathcal G_i$.  Truncate by the substitution in
\cref{eq:binomial-loss}.  We obtain bounded admissible classes
with coefficients $M_b^{\pm}$, for every fixed $b$, satisfying
Kummer conditions at both $p$-places.  The split Heegner norm
relation gives augmentation equal to the Kummer class of
\[
 (a_{q_i}-2)y.
\]
Indeed the $q_i+1$ cyclic subgroups in the Hecke correspondence
comprise two horizontal CM isogenies, corresponding to the two
prime ideals over $q_i$, and the $q_i-1$ descending isogenies
whose target order has conductor $q_i$.  The latter are the
orbit over the conductor-one ring class field.  Thus the relative
trace is the Hecke action minus the two horizontal translates.
Tracing further to $K$ turns each horizontal translate into $y$,
and the Hecke eigenvalue on $A_0$ is $a_{q_i}$.  This proves the
displayed relation; see also the correspondence and CM action
conventions in \cite[Chapter~I, Sections~2--3]{GrossZagier1986}.
Since $y$ is a fixed torsion point, one
fixed scalar kills all these augmentations in integral Kummer
cohomology.  Their rational augmentation is zero.

For $b_*=1$ this proves the cohomological assertion.  In the
strict case divide the limiting Kummer trace once by $t$, as in
\cref{lem:admissible-specialization}; its first specialization,
say $z_1$, is an ordinary compact Selmer class of $V/K$.  The
cohomology sequence of \cref{eq:coefficient-sequence} and
$H^0(K,V\otimes k)=0$ make multiplication by $t$ injective on
$H^1(M_{b-1}^{\pm})\to H^1(M_b^{\pm})$, so this divided class
is unique.  All local Kummer conditions survive this division because the tame
character is trivial at $p$.  The operation merely shifts bounded
coordinates after subtraction of a fixed bounded coboundary; no
denominator depending on $[H_i:K]$ is introduced.

We verify the conjugation identity with its effect on this
coefficient, using the CM coordinates of
\cite[Chapter~II, Section~1]{GrossZagier1986}.  Write a CM point
in fractional-ideal coordinates as
\[
 x(\mathfrak a,\mathfrak n)=
  (\C/\mathfrak a\longrightarrow\C/\mathfrak n^{-1}\mathfrak a),
 \qquad \mathfrak n\overline{\mathfrak n}=(N).
\]
The dual cyclic isogeny gives
$w_Nx(\mathfrak a,\mathfrak n)
 =x(\mathfrak n^{-1}\mathfrak a,\overline{\mathfrak n})$.
Consequently
\[
 x(\mathfrak a,\mathfrak n)^c
   =w_Nx(\mathfrak n\overline{\mathfrak a},\mathfrak n).
\]
In the ring class group,
$[\overline{\mathfrak a}]=[\mathfrak a]^{-1}$, so this is inversion
followed by a fixed translation.  The level ideal is invertible
in every order in question because $(q_i,N)=1$.  If $\varphi$ is
the modular parametrization, the differentials of
$\varphi\circ w_N$ and
$-\operatorname{sign}(A_0)\varphi$ agree.  Their difference is
therefore a constant $Q$, the image of a cusp.  The Manin--Drinfeld
theorem \cite{Manin1972,Drinfeld1973} makes
$Q$ of fixed torsion order.  In the strict case $r=1$, so the
Fricke scalar $-\operatorname{sign}(A_0)$ is $+1$.
Conjugation also inverts the anticyclotomic weight, namely
\begin{equation}\label{eq:weight-involution}
 t\longmapsto(1+t)^{-1}-1=-t+t^2-\cdots.
\end{equation}
The translation is multiplication by a group-ring unit of
augmentation one.  Because the constant Kummer coefficient has
already vanished, that unit has no effect on the leading linear
coefficient.  Equation~\eqref{eq:weight-involution} thus gives
$z_1^c=-z_1$.  But the compact Selmer line is generated by the
invariant class $a$.  Therefore $z_1=0$.

Divide once more by $t$.  Its specialization $z_2$ again belongs
to the compact Selmer line, by the same local Kummer and bounded
division argument.  That line is everywhere locally zero in
the strict case.  Thus the original trace has zero local Kummer
class through degree two, which is the assertion required for
$b_*=3$.

It remains to convert the local Kummer vanishing into
\cref{eq:heegner-log}.  At $p$,
Proposition~\ref{prop:tame-primes} makes $H_i/K$ unramified with
local degree of bounded $p$-part, so
\cref{lem:uniform-local-log} applies with a common constant $c$.
Restricting the local twisted Kummer trace to the unramified field
of definition makes the coefficient action trivial.  It becomes
exactly the Kummer class of the raw weighted sum, modulo the
precision allowed by \cref{eq:binomial-loss}.  The vanishing
through degree $b_*-1$ holds to precision tending to infinity as
in \cref{eq:precision-choice}.  After one fixed factor $p^C$
clears the cochain denominators, the local comparison gives the
lower bound $n-C-c$ at Kummer precision $n$.  Thus the
corresponding logarithm coefficients have valuations tending to
infinity.  For any fixed higher truncation, the Kummer cocycles,
logarithms, and binomial coefficients all have a fixed denominator,
so those coefficients
remain bounded as well.  This proves the proposition.
\end{proof}

When a subsequent construction squares a series, we form the
preceding traces first at any fixed order $b\ge2b_*$.  Boundedness
of all the retained coefficients and
\cref{eq:heegner-log} then permit the required product estimates
without an interchange of an unbounded truncation and a limit.
